\subsubsection{2.8. Thang đo hiệu năng cho dữ liệu mất cân bằng}
\label{sec:metrics}

\indent\indent Trong bài toán hiện tại, dữ liệu thực tế luôn tồn tại sự chênh lệch rất lớn. Do bài toán được thiết lập dưới dạng phân loại đa lớp với $C=4$ nhãn, việc chỉ dựa vào tỷ lệ dự đoán đúng là không đủ để đánh giá toàn diện. Hệ thống đánh giá trong nghiên cứu này được xây dựng dựa trên Ma trận nhầm lẫn đa lớp (\textit{Multi-class Confusion Matrix}) \cite{fawcett2006introduction}. Kế thừa hệ thống thang đo của Sokolova và Lapalme \cite{sokolova2009systematic}, đề tài thiết lập các tiêu chuẩn đánh giá theo chiến lược "Một-chống-lại-tất-cả" (\textit{One-vs-Rest - OvR}) cho từng lớp $i$ như sau:\vspace{-0.1cm}

\begin{itemize}[label={--}, itemsep=1pt]
    \item \textbf{Accuracy (Độ chính xác tổng thể):} Là tỷ lệ tổng số mẫu được dự đoán đúng trên toàn bộ tập dữ liệu. Trong bối cảnh phân bố lớp bất đối xứng, chỉ số này không mang tính quyết định vì mô hình chỉ cần thiên vị hoàn toàn lớp đa số (\texttt{BENIGN}) là đã có thể đạt điểm số rất cao \cite{sokolova2009systematic}. Do đó, Accuracy chỉ được sử dụng làm tham chiếu bổ trợ.
    \[
        \text{Accuracy} = \frac{\sum_{i=1}^{C} TP_i}{N} \quad (\text{Với } N \text{ là tổng số mẫu})
    \]

    \item \textbf{Precision (Độ chuẩn xác) của lớp $i$:} Trả lời câu hỏi: \textit{"Trong số các tài khoản bị hệ thống dự đoán thuộc nhãn $i$, có bao nhiêu phần trăm thực sự mang nhãn $i$?"}. Đặc biệt với các lớp rủi ro cao, nếu Precision thấp, hệ thống đang bắt oan (\textit{False Positive}) người dùng thực tế, gây ảnh hưởng nghiêm trọng.
    \[
        \text{Precision}_i = \frac{TP_i}{TP_i + FP_i}
    \]

    \item \textbf{Recall (Độ phủ) của lớp $i$:} Trả lời câu hỏi: \textit{"Trong tổng số tài khoản thực sự thuộc nhãn $i$, hệ thống đã nhận diện đúng được bao nhiêu phần trăm?"}. Đối với các nhóm \textit{bot farm} (\texttt{MALICIOUS}), nếu Recall thấp, hệ thống bị lọt lưới (\textit{False Negative}) nhiều.
    \[
        \text{Recall}_i = \frac{TP_i}{TP_i + FN_i}
    \]

    \item \textbf{Macro F1-Score:} Giữa Precision và Recall luôn tồn tại sự đánh đổi tự nhiên. Tương tự các nghiên cứu phát hiện Sybil đương đại \cite{liu2025detecting, heeb2024sybil}, đề tài áp dụng Macro F1-Score để tính trung bình cộng điểm F1 của toàn bộ $C$ lớp. Chỉ số này đối xử bình đẳng với mọi lớp bất chấp số lượng mẫu, ép mô hình phải hoạt động xuất sắc ở cả nhiệm vụ nhận diện được cụm Sybil lẫn bảo vệ người dùng thật.
    \[
        \text{F1}_i = 2 \times \frac{\text{Precision}_i \times \text{Recall}_i}{\text{Precision}_i + \text{Recall}_i} \quad ; \quad \text{Macro F1} = \frac{1}{C} \sum_{i=1}^{C} \text{F1}_i
    \]

    \item \textbf{ROC AUC Đa lớp (Multi-class AUC):} Vì đầu ra của GNN là phân phối xác suất trên 4 lớp (thông qua hàm \textit{Softmax}), việc đánh giá dựa trên một ngưỡng cứng mang tính chủ quan. Theo chứng minh của Bradley \cite{bradley1997use}, diện tích dưới đường cong ROC (AUC) đo lường khả năng phân tách của mô hình một cách khách quan. Đối với bài toán đa lớp, đề tài sử dụng chiến lược \textit{Macro-AUC (OvR)}, tính toán khả năng mô hình phân biệt một lớp cụ thể với 3 lớp còn lại, đảm bảo tính bền bỉ trước sự mất cân bằng dữ liệu.
\end{itemize}